\section{Qualit\`a delle mesh anatomiche e discrepanza regionale}\label{app:anatomy}
\subsection*{Interfacce geometriche e classificazione dei materiali}
La segmentazione BraTS utilizza le etichette $0,1,2,4$ per identificare background del dominio, NCR/NET, edema ed enhancing tumor. I moltiplicatori della Tabella~\ref{tab:materials} sono assegnati alle celle quadrilatere secondo la rispettiva etichetta; le interfacce tra materiali sono pertanto approssimate da facce del reticolo e non imposte come frontiere esatte della segmentazione voxel. Le mesh anatomiche dell'inversione hanno passo obiettivo $3\,$mm e $2\,$mm e comprendono rispettivamente $1901$ e $4287$ quadrilateri.

L'audit geometrico compara le regioni come insiemi poligonali, dopo aver ricostruito le celle dai file Gmsh delle due discretizzazioni anatomiche. La sovrapposizione del dominio esterno \`e molto alta (Dice $0.99795$), ma la coincidenza delle singole regioni materiali risulta inferiore, come mostra la Tabella~\ref{tab:app-mesh-compare}.
\begin{table}[H]
\centering\caption{Confronto delle regioni materiali delle due mesh anatomiche del confronto inverso; Dice calcolato tra insiemi poligonali con identica etichetta, non rispetto alla segmentazione voxel originale.}\label{tab:app-mesh-compare}
\small\setlength{\tabcolsep}{4pt}
\begin{tabular}{@{}lrrr@{}}\toprule
Regione&Celle $3\,$mm&Celle $2\,$mm&Dice\\\midrule
Background&1339&3049&0.97783\\
NCR/NET&66&153&0.85406\\
Edema&399&875&0.91268\\
ET&97&210&0.79977\\\bottomrule
\end{tabular}
\end{table}

Sull'area comune alle due discretizzazioni, circa $4.71\%$ della superficie riceve etichette differenti. Questo fatto determina una variazione del campo dei coefficienti per salti alle interfacce, anche quando la configurazione nominale e i moltiplicatori regionali sono invariati. Il confronto spaziale non rappresenta pertanto una stima pura dell'errore FEM su un campo di coefficienti immutabile.

\subsection*{Sensori congelati e localizzazione della discrepanza}
Per il protocollo bilanciato, tre sensori su sedici cambiano classe regionale fra le mesh. Due posizioni spiegano circa l'$86.5\%$ della somma dei quadrati della differenza fra i vettori osservativi del caso $3\,$mm/$48$ e del caso $2\,$mm/$96$. Nel protocollo focalizzato nell'edema nessun sensore cambia etichetta, ma due posizioni rappresentano circa il $95\%$ della stessa discrepanza quadratica. La sola appartenenza regionale del punto di osservazione non determina dunque l'intera sensibilit\`a alla mesh: incidono anche gradienti vicini, posizione delle interfacce e approssimazione dei campi. I sensori vengono mantenuti congelati e non eliminati dopo aver esaminato gli errori.

Le differenze relative osservazionali sono riassunte nella Tabella~\ref{tab:app-discrepancies}. I confronti temporali alla stessa mesh sono entrambi intorno allo $0.02\%$, mentre il confronto con variazione spaziale \`e dell'ordine del $2\%$ per i sensori bilanciati. Le quattro traiettorie di riferimento presentano valori nodali di concentrazione finiti e non negativi, con deriva relativa della massa pesata della concentrazione dell'ordine di $10^{-15}$.
\begin{table}[H]
\centering\caption{Differenze relative tra vettori osservativi del modello M4 anatomico eterogeneo, protocollo bilanciato; dati sintetici senza rumore.}\label{tab:app-discrepancies}
\small\setlength{\tabcolsep}{4pt}
\begin{tabular}{@{}lr@{}}\toprule
Confronto&$\|y_a-y_b\|_2/\|y_b\|_2$\\\midrule
$3\,$mm/$48$ vs $3\,$mm/$96$&$0.019\%$\\
$2\,$mm/$48$ vs $2\,$mm/$96$&$0.020\%$\\
$3\,$mm/$48$ vs $2\,$mm/$96$&$2.035\%$\\\bottomrule
\end{tabular}
\end{table}

\subsection*{Decomposizione dell'errore spaziale}
Indichiamo con $m_{h}$ il campo dei materiali approssimato per etichette di cella. La differenza osservata
\begin{equation}\label{eq:crossmesh-mixed}
\delta y=\Gmap_{h_2}(\tht^\dagger;m_{h_2})-
\Gmap_{h_3}(\tht^\dagger;m_{h_3})
\end{equation}
mescola la dipendenza da $h$ e quella da $m_h$. Per separarle occorre definire \emph{una sola} funzione materiale fisica $m^*(\mathbf x)$, per esempio valutando la medesima segmentazione voxel ai punti di quadratura sulle differenti mesh. A quel punto la componente numerica sarebbe definita come
\begin{equation}\label{eq:crossmesh-clean}
\delta y_{\rm FEM}=\Gmap_{h_2}(\tht^\dagger;m^*)-
\Gmap_{h_3}(\tht^\dagger;m^*).
\end{equation}
Un terzo livello spaziale con $m^*$ fissato consentirebbe anche uno studio dell'ordine empirico e dell'errore osservativo, purch\'e siano controllate regolarit\`a e approssimazione del contorno esterno. Le quantit\`a della \eqref{eq:crossmesh-clean} non sono disponibili nei test eseguiti: l'analisi quantitativa del testo riguarda la \eqref{eq:crossmesh-mixed}, che rimane un controllo valido della discrepanza totale e della stabilit\`a delle stime rispetto al cambio di rappresentazione anatomica.
